\documentclass[12pt]{article}
\usepackage[utf8]{inputenc}
\usepackage[T1]{fontenc}
\usepackage{hyperref}
\title{Chung Language Resources for K-12 Curriculum Development}
\author{Compiled by Claude Sonnet 5.5 (claude-sonnet-5-5) for Andr\'as Kornai}
\date{}
\begin{document}
\maketitle

\section{Introduction}
Chung (Cung, Chungmboko; ISO 639-3 cnq) is a Beboid language of Cameroon. Speaker estimates differ: 1,400 (2001, Wikipedia), 2,620 (Global Recordings Network), 2,700 (Joshua Project). Very little material exists; no school curriculum resources were found.

\section{Primary Educational Resources}
None were found.

\section{Linguistic and Reference Materials}
\begin{itemize}
\item Wikipedia overview: \url{https://en.wikipedia.org/wiki/Chung_language}
\item Joshua Project people-group profile: \url{https://joshuaproject.net/people_groups/22322}
\end{itemize}

\section{Religious and Community Resources}
\begin{itemize}
\item Scripture Earth lists a Chung New Testament: \url{https://www.scriptureearth.org/00eng.php?iso=cnq}
\end{itemize}

\section{Digital and Audio Resources}
\begin{itemize}
\item Global Recordings Network audio recordings: \url{https://globalrecordings.net/en/language/cnq}
\end{itemize}

\section{Curriculum Development Context}
The language is spoken by a small community and is not known to be taught in schools. A curriculum would need new primer and orthography development, building on the New Testament text as the main written corpus.
\end{document}